% David Koch

\documentclass[
	headings=optiontotocandhead,% Erweiterung für das optionale Argument der
	% Gliederungsbefehle aktiviert.
	oneside,
	numbers=noenddot,% Keine Punkte am Ende der Gliederungsnummern und davon
	% abgeleiteten Nummern
	toc=flat, %Flache TOC --- kann man anpassen (auskommentieren)
	10pt, % Schriftgröße
	parskip=full, % Abstand zwischen Absätzen (ganze Zeile)
	listof=totoc, % Verzeichnisse im Inhaltsverzeichnis aufführen
	listof=flat, % mehr Abstand für grosse Zahlen
	numbers=noenddot, % kein Punkt am Ende bei Nummern
	%%enlargefirstpage,% Gibt es bei scrartcl nicht!!!!
	bibliography=totoc, % Literaturverzeichnis im Inhaltsverzeichnis aufführen
	%index=totoc, % Index im Inhaltsverzeichnis aufführen
	%captions=tableheading, % Beschriftung von Tabellen für Ausgabe oberhalb
	% der Tabelle formatieren
	%draft % Status des Dokuments (final/draft) draft hinzufügen zum anziegen
	%%der zeilen ende
	a4paper,DIV=14,
	% captions=tablesignature,
]{scrartcl}

\setcounter{secnumdepth}{3}

\usepackage[T1]{fontenc}
\usepackage[utf8]{inputenc}

\usepackage[english, ngerman]{babel, varioref} % your native language must be the last one!!

\usepackage{lastpage}
\usepackage{listings}
\usepackage{blindtext}

%% Aufzählungen nicht so weit einrücken
\usepackage[inline]{enumitem}
%\setitemize{leftmargin=*}
% Listen etwas wenige einrücken, erfordert enumitem
\setitemize{labelindent=2em,labelsep=0.5cm,leftmargin=12ex}

\usepackage{lmodern}

\usepackage{xspace}

\usepackage{graphicx}
\graphicspath{ {../assets} }

%%? \usepackage{textcomp}
\usepackage[hyphens]{url}
\usepackage{makeidx}
\makeindex
%%? \usepackage{graphicx}
\usepackage[numbers]{natbib}
\PassOptionsToPackage{normalem}{ulem}
\usepackage{ulem}

\usepackage{needspace}

\setlength\partopsep{0.5ex}%schoenere Listen
\usepackage[bottom]{footmisc}%fussnote ganz unten

\usepackage[]{microtype}
\UseMicrotypeSet[protrusion]{basicmath} % disable protrusion for tt fonts

\usepackage{multirow}   % Allows table elements to span several rows.
\usepackage{booktabs}   % Improves the typesettings of tables.
\usepackage{subcaption} % Allows the use of subfigures and enables their referencing.
\usepackage[ruled,linesnumbered]{algorithm2e} % Enables the writing of pseudo code.
\usepackage[usenames,dvipsnames,table]{xcolor} % Allows the definition and use of colors. This package has to be included before tikz.
\usepackage{nag}       % Issues warnings when best practices in writing LaTeX documents are violated.
\usepackage{todonotes} % Provides tooltip-like todo notes.

\usepackage{color}
\usepackage[binary-units]{siunitx}

%% Override default figure placement To be within the flow of the text rather
%% than on it's own page.
% \usepackage{float}
% \makeatletter
% \def\fps@figure{H}
% \makeatother

%% bei vielen Bildern o.ä sinnvoll: Seite muss nicht bis ganz unten gefüllt werden
% \raggedbottom

%\usepackage{footbib} %  footcite, needs other tooling
%% for pandoc2 images
\makeatletter
\def\maxwidth{\ifdim\Gin@nat@width>\linewidth\linewidth\else\Gin@nat@width\fi}
\def\maxheight{\ifdim\Gin@nat@height>\textheight\textheight\else\Gin@nat@height\fi}
\makeatother
% Scale images if necessary, so that they will not overflow the page
% margins by default, and it is still possible to overwrite the defaults
% using explicit options in \includegraphics[width, height, ...]{}
\setkeys{Gin}{width=\maxwidth,height=\maxheight,keepaspectratio}

%% bessere Suche im PDF
\input{glyphtounicode}
\pdfgentounicode=1
%%%%%%%%%%%%%%%%%%%%%%%%%%%%%%%%%%%%%%%%%%%%%%%%%%%%%%%%%%%%%%%%%%%%%%%%%%%%%%%%%%

%  Kopf und Fußzeilen -- links und rechts verschieden
\newcommand{\kopfbild}{\voffset7mm\includegraphics[width=25mm]{HTL3RLogo}}
\newcommand{\kopfHTL}{\sffamily{\textbf{\large{Projekthandbuch HTL3R}}}}

\usepackage[automark,footsepline,plainfootsepline]{scrlayer-scrpage}
\setkomafont{pageheadfoot}{\normalcolor\footnotesize\scshape}
\setkomafont{pagenumber}{\normalfont\normalsize}
\clearpairofpagestyles
\ihead{\headmark}
\ihead{\kopfbild}
\ohead{\kopfHTL}
\ifoot{\smaller{Höhere Technische Bundeslehranstalt Wien 3 | Rennweg 89b | 1030 Wien | \textcolor{orange}{www.htl.rennweg.at}}}
\ofoot{Seite \pagemark/\pageref{LastPage}}
\ModifyLayer[addvoffset=-.6ex]{scrheadings.foot.above.line}% Linie verschieben
\ModifyLayer[addvoffset=-.6ex]{plain.scrheadings.foot.above.line}% Linie verschieben
\setlength{\headheight}{32pt}

% alle Seiten mit Kopfzeile
%\renewcommand{\chapterpagestyle}{scrheadings}

%% Code Beispiele
%% eine Variante
\usepackage{listings}
\renewcommand{\lstlistingname}{\inputencoding{utf8}Listing}

\usepackage{tabularx}
\usepackage{scrhack}

\usepackage{array}
\newcommand\Tstrut{\rule{0pt}{3.2ex}}         % = `top' strut
\newcommand\Bstrut{\rule[-1.5ex]{0pt}{0pt}}   % = `bottom' strut

\newenvironment{nstabbing}
	{\setlength{\topsep}{-\parskip}
		\setlength{\partopsep}{-\parskip}
		\tabbing}
	{\endtabbing}

\usepackage{titlesec}
% \titleformat{?Überschriftenklasse?}[Absatzformatierung?]{?Textformatierung?} {?Nummerierung?}{?Abstand zwischen Nummerierung und Überschriftentext?}{?Code vor der Überschrift?}[?Code nach der Überschrift?]
\titleformat{\section}[hang]{\Large\bfseries\sffamily}{\thesection\quad}{-1.2ex}{}
\titleformat{\subsection}[hang]{\large\bfseries\sffamily}{\thesubsection\quad}{-1.2ex}{}
\titleformat{\subsubsection}[hang]{\large\bfseries\sffamily}{\thesubsubsection\quad}{-1.2ex}{}
\titleformat{\paragraph}[hang]{\large\bfseries\sffamily}{\theparagraph\quad}{-1.2ex}{}

% \titlespacing{?Überschriftenklasse?}{?Linker Einzug?}{?Platz oberhalb?}{?Platz unterhalb?}[?rechter Einzug?]
\titlespacing{\section}{0pt}{6pt}{6pt}
\titlespacing{\subsection}{0pt}{6pt}{0pt}
\titlespacing{\subsubsection}{0pt}{6pt}{0pt}
\titlespacing{\paragraph}{0pt}{6pt}{0pt}

%% sollte das letzte Package sein
\usepackage[unicode=true,
bookmarks=true,bookmarksnumbered=false,bookmarksopen=false,
breaklinks=true,pdfborder={0 0 0},backref=false,colorlinks=false]
{hyperref}
\hypersetup{pdftitle={Fenrir Management Review vom 10.09.24},
	pdfauthor={Koch, Uhlig, Burger, Vogler},
	pdfsubject={Diplomarbeit},
	pdfkeywords={5CN, DA, Fenrir}}
\urlstyle{same} % don't use monospace font for urls

% Auch Fußnoten bündig ausrichten
\deffootnote[]{1em}{1em}{\textsuperscript{\thefootnotemark\ }}
%% setup
\sloppy % weniger Meldungen
\voffset7mm % etwas nach unten

%%%%%%%%%%%%%%%%%%%%%%%%%%%%%%%%%%%%%%%%%%%%%%%%%%%%%%%%%%%%%%%%%%%%%%%%%%%%%%%%%%
\begin{document}
%% schöner: 10000 -- gar keine, 1000 als Mittelweg
\clubpenalty = 10000 % Schusterjungen verhindern
\widowpenalty = 10000 % Hurenkinder verhindern
\displaywidowpenalty = 1000

{\sffamily{\textbf{\LARGE{\textcolor{orange}{Management Review}}}}}\\
\noindent\rule{\textwidth}{0.1pt}
\begin{nstabbing}
	\hspace{4cm}\=\hspace{4cm}\=\hspace{4cm}\=\kill
	Projekttitel: \> \textbf{Fenrir}\\
	Auftraggeber: \> \textbf{Christian Schöndorfer}\\
	Auftragnehmer: \> \textbf{David Koch}\\
	Schuljahr: \> \textbf{2024/25}
	\> Klasse: \> \textbf{5CN}\\
\end{nstabbing}
{\smaller
	\begin{tabularx}{\textwidth}{l l l l}
	\hline
	\textbf{Version} & \textbf{Datum} & \textbf{Autorin/Autor} & \textbf{Änderung}\Tstrut  \\
	v1.0 & 09.09.2024 & David Koch & Erstellung des Management Reviews (+ Vorlage)\Bstrut \\
	\hline
	\end{tabularx}
}

\section{Projektstatus}
\textbf{Berichtszeitraum: 29.06.2024 bis 10.09.2024}

Kurze verbale Beschreibung, was wurde erreicht, welche Meilensteine, was nicht, welche Probleme. Wo steht das Projekt; 

Derzeitiger Ampelstatus: Grün

\textbf{Termine}

{\smaller
	\begin{tabularx}{\textwidth}{|X|X|X|X|X|}
		\hline
		\textbf{\% abgeschlossen} & \textbf{Projektstatus} & \textbf{Erwartetes\newline Projektende} & \textbf{Geplantes Projektende} & \textbf{Abweichung} \\
		\hline
		- & - & - & - & - \\
		\hline
	\end{tabularx}
}
% \% abgeschlossen = Prozentwert, wie viel des Gesamtprojektes abgeschlossen ist. Ermittlung auf der Basis der abgeschlossenen Arbeiten (Stunden) im Verhältnis zu den noch zu erledigenden Arbeiten (in Stunden)

\textbf{Kosten}

{\smaller
	\begin{tabularx}{\textwidth}{|X|X|X|X|X|}
		\hline
		\textbf{Ist-Kosten} & \textbf{Restkosten} & \textbf{Erwartete\newline Gesamtkosten} & \textbf{Plankosten} & \textbf{Abweichung} \\
		\hline
		- & - & - & - & - \\
		\hline
	\end{tabularx}
}
% Ist-Kosten = bisher angefallene Kosten in € / bisher angefallene Stunden \\
% Restkosten = welche Kosten bzw. wieviel Stunden sind ab heute bis zum Projektende noch zu erwarten \\
% Erwartete Gesamtkosten = Ist-Kosten + Restkosten bzw. Ist-Aufwand (h) + Restaufwand (h) \\
% Plankosten = ursprünglich (Projektbeginn, Angebot) berechnete Kosten bzw. ursprünglich geschätzter Stundenaufwand \\
% Abweichung = Erwartete Gesamtkosten – Plankosten in € bzw. Erwarteter Gesamtaufwand – Planaufwand (h) \\

\textbf{Stunden}

{\smaller
	\begin{tabularx}{\textwidth}{|X|X|X|X|X|}
		\hline
		\textbf{Ist-Aufwand} & \textbf{Restaufwand} & \textbf{Erwarteter\newline Gesamtaufwand} & \textbf{Planaufwand} & \textbf{Abweichung} \\
		\hline
		- & - & - & - & - \\
		\hline
	\end{tabularx}
}
% Analog zu den Kosten (siehe oben)

\subsection{Teammotivation \colorbox{green!30}{:)}} 
Die Motivation ist im gesamte Team hoch, es bestehen jedoch Sorgen, dass durch zunehmenden Stress im Schuljahr die Arbeitszeit für die Diplomarbeit zu sehr beeinträchtigt werden kann.

\section{Probleme im Projekt}
Beim Aufbau der Prototyp-Betriebszelle sind der Wasserfüllstandssensor sowie der Raspberry Pi beschädigt worden (Totalschaden). Die Ersatzbauteile wurden bereits bestellt und verbaut, somit besteht keine Beeinträchtigung des Projektablaufs, es bleiben mit diesen Schäden jedoch mehr Kosten offen.

Hier erzählt Burger über die Wehen mit vSphere und terraform falls sie schlimm genug waren AAAAAAAAAa

Gabi bitte hier über 3D-Drucker beschweren AAAAAAa

\subsection{Problemlösungsstrategie}
Alle im Projekt bisher aufgetretenen Probleme wurden bereits gelöst, somit ist keine Problemlösung notwendig.

\subsection{Handlungsbedarf seitens des Managements}
Sind Aktivitäten seitens des Managements (Projektauftraggeber, Projektbetreuung, Geschäftsleitung, Management) notwendig? Wenn ja, welche?

\begin{enumerate}
	\item Abnahme des Pflichtenhefts bis .......
	\item Entscheidung über Auswahl des Screendesigns bis .......
\end{enumerate}

\section{Erledigte Arbeiten (vollständig)}
\begin{table}[h]
	\begin{tabularx} {\textwidth} {
			|>{\hsize=.12\hsize}X
			|>{\hsize=.08\hsize}X
			|>{\hsize=.34\hsize}X
			|>{\hsize=.10\hsize}X
			|>{\hsize=.12\hsize}X
			|>{\hsize=.12\hsize}X
			|>{\hsize=.09\hsize}X|
		}
		
		\hline
		\rowcolor[HTML]{D9D9D9} 
		\textbf{\normalsize{Bearbeiter}} & \textbf{\normalsize{PSP-Code}} & {\textbf{\normalsize{Tätigkeit}}} & \textbf{\normalsize{Ort}} & \textbf{\normalsize{Dauer geplant (h)}} & \textbf{\normalsize{Dauer benötigt (h)}} & \textbf{\normalsize{Status}} \\ \hline
		- & - & - & - & - & - & \cellcolor{green!30} \\ \hline
	\end{tabularx}
\end{table}

Die während den Sommerferien erledigten Arbeiten bestehen hauptsächlich aus dem Aufbau der Prototyp-Betriebszelle (inklusive dem Anlernen der Grundlagen bezüglich der Programmierung und Steuerung von Bussystemen wie OneWire, I2C und Modbus), der Einrichtung von vSphere-Provisionierung mittels Terraform auf der UCS (in Kooperation mit dem Diplomarbeitsteam Atropos) sowie der Erstellung der Diplomarbeits-Website. Diese ist bereits unter der URL \url{https://www.fenrir-ot.at/} erreichbar.



\section{Anstehende Arbeiten}
\begin{table}[h]
	\begin{tabularx} {\textwidth} {
			|>{\hsize=.12\hsize}X
			|>{\hsize=.08\hsize}X
			|>{\hsize=.54\hsize}X
			|>{\hsize=.12\hsize}X
			|>{\hsize=.15\hsize}X|
		}
		
		\hline
		\rowcolor[HTML]{D9D9D9} 
		\textbf{\normalsize{Bearbeiter}} & \textbf{\normalsize{PSP-Code}} & {\textbf{\normalsize{Tätigkeit}}} & \textbf{\normalsize{Dauer geplant (h)}} & \textbf{\smaller{Fertigstellung geplant}} \\ \hline
		- & - & - & - & - \\ \hline
	\end{tabularx}
\end{table}

Die zuvor erstellten Planungsdokumente müssen mit den über die Sommerferien erhaltenen Informationen bezüglich des Projektumfangs überarbeitet werden. Wie zuvor erwähnt muss bis zum 20.09.2024 der Diplomarbeitsantrag offiziell abgeschickt werden, somit sollten spätestens eine Woche zuvor (am 13.09.2024) beide Diplomarbeitsbetreuer eine beinahe fertige Kopie erhalten, um diese auf Fehler zu überprüfen. Da im Antrag die Zieldefinition im Ganzen enthalten ist, hat die Fertigstellung/Überarbeitung -- im Vergleich zu den restlichen Planungsdokumenten -- dieser die höchste Priorität.

Der Aufbau der nächsten OT-Betriebszelle beginnt, sobald der vom Kollegen Vogler bestellte 3D-Drucker ankommt und eine maßgeschneiderte Förderschnecke drucken kann.

Es muss auch eine dedizierte VPN-Verbindung zum Kollegen Vogler nach Hause erstellt werden (mittels Twingate), um jederzeit auf die bestehende OT-Gerätschaft zugreifen zu können, als auch diese später mit virtuellen Maschinen auf der UCS verbinden zu können.

\end{document}